\begingroup

All schemes and morphisms in this chapter are ordinary schemes and scheme
morphisms. A variety over a field is integral, separated and of finite type;
geometric integrality is not imposed, following \cite{deJong1996}, 2.9.

\begin{definition}[Variety {\cite[p.~55]{deJong}}]
\label{def:variety}
\lean{Alterations.IsVariety}
\leanok
A variety over a field $k$ is an integral scheme equipped with a separated,
finite-type structure morphism to $\operatorname{Spec}k$.
\end{definition}

\endgroup

\begingroup

\begin{lemma}[The generic point misses a proper support]
\label{lem:integral-ideal-generic-outside}
\uses{def:variety}
\lean{Alterations.Geometry.IntegralIdeal.genericPoint_not_mem_support}
\leanok
Let $X$ be integral and let $\mathcal I$ be a quasi-coherent ideal sheaf
with proper closed support. The generic point of $X$ is outside this support.
\end{lemma}
\begin{proof}
If the generic point belonged to this closed set, every specialization
of it would belong too. Every point of the integral scheme is such a
specialization, making the support all of $X$, a contradiction.
\end{proof}

\endgroup

\begingroup

\begin{lemma}[Every nonempty affine coordinate ideal is nonzero {\cite[p.~67]{deJong}}]
\label{lem:integral-ideal-affine-nonzero}
\uses{lem:integral-ideal-generic-outside}
\lean{Alterations.Geometry.IntegralIdeal.ideal_ne_bot_of_support_ne_top}
\leanok
Under the preceding hypotheses, the actual coordinate ideal of
$\mathcal I$ on every nonempty affine open $U$ is nonzero.
\end{lemma}
\begin{proof}
The generic point lies in $U$, since every nonempty open of an integral
scheme contains it. If the coordinate ideal were zero, its vanishing
set on $U$ would be all of $U$. The exact affine description of ideal-sheaf
support would then place the generic point in the support, contradicting
the preceding lemma.
\end{proof}
\endgroup

\begingroup

\begin{lemma}[Reduced projective spectrum of a domain]
\label{lem:proj-domain-reduced}
\lean{Alterations.Geometry.ProjIntegral.isReduced_proj}
\leanok
Let $A$ be a domain graded by the nonnegative integers. Then
$\operatorname{Proj}A$ is reduced, including when it is empty.
\end{lemma}
\begin{proof}
At a homogeneous prime $\mathfrak p$, the local ring of the projective
spectrum is the degree-zero homogeneous localization, a subring of the
ordinary localization $A_{\mathfrak p}$. That localization is a domain,
so its subring is reduced. Reducedness of every stalk implies reducedness
of the scheme.
\end{proof}

\endgroup

\begingroup

\begin{lemma}[Generic point of a projective spectrum]
\label{lem:proj-domain-generic-point}
\lean{Alterations.Geometry.ProjIntegral.zeroPoint,Alterations.Geometry.ProjIntegral.isGenericPoint_of_asHomogeneousIdeal_eq_bot}
\leanok
Let $A$ be a graded domain with nonzero irrelevant ideal. The homogeneous
zero prime is a point of $\operatorname{Proj}A$ and is its generic point.
\end{lemma}
\begin{proof}
The zero ideal is homogeneous and prime. The irrelevant ideal is not
contained in zero because it is nonzero, so this prime is relevant.
The specialization order on the projective spectrum is inclusion of
homogeneous prime ideals. The zero ideal is contained in every such
prime, hence the closure of this point is the whole projective spectrum.
\end{proof}

\endgroup

\begingroup

\begin{theorem}[Integral projective spectra]
\label{thm:proj-domain-integral}
\uses{lem:proj-domain-reduced,lem:proj-domain-generic-point}
\lean{Alterations.Geometry.ProjIntegral.isIntegral_proj}
\leanok
The projective spectrum of a graded domain with nonzero irrelevant ideal
is an integral scheme.
\end{theorem}
\begin{proof}
The generic point makes the underlying space nonempty and irreducible.
The preceding stalk argument makes the scheme reduced. An irreducible
reduced scheme is integral.
\end{proof}

\endgroup

\begingroup

\begin{definition}[The Rees algebra and its degree pieces]
\label{def:affine-rees-grading}
\lean{Alterations.reesGrading}
\leanok
Let $R$ be a commutative ring and $I\subset R$ an ideal. Its Rees algebra is the actual
subalgebra
\[
\mathcal R(I)=\left\{\sum_n a_nT^n\in R[T]\;\middle|\;a_n\in I^n\text{ for every }n\right\}.
\]
Define $\mathcal R(I)_n=\mathcal R(I)\cap RT^n$, so that
$\mathcal R(I)_n=I^nT^n$ with its inherited $R$-module structure.
\end{definition}

\endgroup

\begingroup

\begin{lemma}[The genuine Rees grading]
\label{lem:affine-rees-grading}
\uses{def:affine-rees-grading}
\lean{Alterations.reesGradedAlgebra,Alterations.reesGrading_proj}
\leanok
The degree pieces form an internal graded algebra structure on $\mathcal R(I)$.
The degree-$n$ component of $p$ is exactly $a_nT^n$, where $a_n$ is the coefficient
of $T^n$ in $p$.
\end{lemma}
\begin{proof}
Each component is in the Rees algebra since $a_n\in I^n$ and its other coefficients
vanish. Only finitely many components are nonzero, and their sum is $p$ by equality
of polynomial coefficients. The decomposition is unique: the coefficient of $T^n$
in a sum of homogeneous elements is the coefficient of the element in degree $n$.
The unit has degree zero. If $aT^i$ and $bT^j$ belong to the respective pieces,
their product is $abT^{i+j}$ and $ab\in I^iI^j=I^{i+j}$. This proves both the
direct sum and multiplication requirements.
\end{proof}

\endgroup

\begingroup

\begin{lemma}[The degree-zero ring]
\label{lem:affine-rees-degree-zero}
\uses{lem:affine-rees-grading}
\lean{Alterations.reesDegreeZeroEquiv}
\leanok
The map $R\to\mathcal R(I)_0$ sending $a$ to its constant polynomial is an isomorphism
of rings.
\end{lemma}
\begin{proof}
The equality $I^0=R$ shows that every constant is in the Rees algebra. A homogeneous
element of degree zero has no coefficient except its constant coefficient, so
the map is surjective. Equality of constant coefficients proves injectivity.
Addition, multiplication and the identity agree with those of constant polynomials.
\end{proof}

\endgroup

\begingroup

\begin{definition}[Affine blowup]
\label{def:affine-blowup}
\uses{lem:affine-rees-grading,lem:affine-rees-degree-zero}
\lean{Alterations.affineBlowup,Alterations.affineBlowupStructure}
\leanok
The affine blowup of $I$ is the scheme
$\operatorname{Bl}_I(\operatorname{Spec}R)=\operatorname{Proj}\mathcal R(I)$ with its
structure map to $\operatorname{Spec}R$ obtained by identifying the degree-zero ring with $R$.
\end{definition}

\endgroup

\begingroup

\begin{corollary}[Integral affine blowups]
\label{cor:affine-blowup-integral}
\uses{thm:proj-domain-integral,def:affine-blowup}
\lean{Alterations.Geometry.ProjIntegral.isIntegral_affineBlowup}
\leanok
If $R$ is a domain and $I\ne0$ is an ideal, its native affine blowup
$\operatorname{Proj}\bigoplus_{n\ge0} I^nT^n$ is integral.
No finite-generation hypothesis is needed for this assertion.
\end{corollary}
\begin{proof}
The Rees algebra is a subring of the polynomial domain $R[T]$, so is a
domain. A nonzero element $a\in I$ gives a nonzero homogeneous element
$aT$ of degree one in the Rees algebra. Thus its irrelevant ideal is
nonzero. Apply the integral-projective-spectrum theorem.
\end{proof}
\endgroup

\begingroup

Let $X$ be an integral scheme, let $D$ be native ideal-sheaf
data on $X$, and suppose that its closed support is not all
of $X$. Write $B=\operatorname{Bl}_D(X)$ for the actual
relative gluing of the affine Rees projective spectra and
$b:B\to X$ for its structure morphism. Finite generation is
not assumed.

\begin{lemma}[Integral native affine blowup pieces]
\label{lem:global-blowup-integral-affine-pieces}
\uses{lem:integral-ideal-affine-nonzero,cor:affine-blowup-integral}
\lean{Alterations.Geometry.BlowupIntegral.chart_isIntegral}
\leanok
For every nonempty affine open $U\subset X$, the actual
scheme $\operatorname{Bl}_{D(U)}(\operatorname{Spec}\mathcal O_X(U))$
is integral.
\end{lemma}
\begin{proof}
Integrality of $X$ makes $\mathcal O_X(U)$ a domain when
$U$ is nonempty. The proper support hypothesis implies
$D(U)\ne0$. The constructed native Rees blowup
of a nonzero ideal in a domain is integral. Applying that
result gives exactly the affine piece used in the gluing.
\end{proof}

\endgroup

\begingroup

\begin{definition}[The fraction chart {\cite[p.~64, p.~65]{deJong}}]
\label{def:affine-blowup-fraction-chart}
\lean{Alterations.affineBlowupFractionChart}
\leanok
Let $R$ be a domain with fraction field $K$, let $I\subset R$ be an ideal and
let $0\ne g\in I$. Define $R[I/g]\subset K$ to be the $R$-subalgebra generated by
$a/g$ for all $a\in I$.
\end{definition}

\endgroup

\begingroup

The blowup used in Section~4.8 of \cite{deJong1996} is a modification.
The argument below proves dominance directly for the actual affine blowup
of a nonzero ideal in a domain, and proves surjectivity when the ideal is
finitely generated.

\begin{lemma}[The fraction-chart coefficient map is injective]
\label{lem:affine-blowup-chart-coefficients-injective}
\uses{def:affine-blowup-fraction-chart}
\lean{Alterations.affineBlowupFractionChart_algebraMap_injective}
\leanok
For a domain $R$ and any $g\in R$, the coefficient map $R\to R[I/g]$ is
injective.
\end{lemma}
\begin{proof}
The composition with the inclusion $R[I/g]\hookrightarrow\operatorname{Frac}(R)$
is the original injective coefficient map $R\to\operatorname{Frac}(R)$.
Equality of two coefficients in the subalgebra therefore implies equality
of their images in the fraction field, hence equality in $R$.
\end{proof}

\endgroup

\begingroup

\begin{lemma}[Evaluation of the homogeneous chart]
\label{lem:affine-blowup-chart-evaluation}
\uses{lem:affine-rees-grading,def:affine-blowup-fraction-chart}
\lean{Alterations.reesChartEvaluation,Alterations.reesChartEvaluation_mk,Alterations.reesChartEvaluation_injective}
\leanok
The polynomial evaluation $T\mapsto1/g$ induces an injective $R$-algebra map
\[
\mathcal R(I)_{(gT)}\longrightarrow K,
\qquad \frac{aT^n}{(gT)^n}\longmapsto\frac a{g^n}.
\]
\end{lemma}
\begin{proof}
Evaluation is an $R$-algebra map on the Rees algebra, and sends $gT$ to $1$.
It therefore extends to the ordinary localization at $gT$. Restrict this map
to the homogeneous degree-zero localization. A homogeneous numerator in degree
$n$ is uniquely $aT^n$, so its image is the displayed fraction. If that image
is zero, the facts that $g\ne0$ and $R\hookrightarrow K$ is injective imply
$a=0$. The original numerator and consequently its homogeneous fraction are
zero. Applying this zero-kernel argument to the difference of two elements
proves injectivity.
\end{proof}

\endgroup

\begingroup

\begin{lemma}[Fractions of powers of the ideal]
\label{lem:affine-blowup-ideal-power-fraction}
\uses{def:affine-blowup-fraction-chart}
\lean{Alterations.idealPower_div_pow_mem_fractionChart}
\leanok
If $a\in I^n$, then $a/g^n\in R[I/g]$.
\end{lemma}
\begin{proof}
For $n=0$ this is the inclusion of $R$ in the algebra. Suppose the assertion
holds in degree $n$. Every element of $I^{n+1}=I I^n$ is a finite sum of products
$bc$, with $b\in I$ and $c\in I^n$. The fraction associated with such a product
is $(b/g)(c/g^n)$, which belongs to the algebra by its definition and the
induction hypothesis. Addition of the products proves the assertion for the
given element of $I^{n+1}$.
\end{proof}

\endgroup

\begingroup

\begin{lemma}[The homogeneous coordinate ring is the fraction chart]
\label{lem:affine-blowup-chart-ring}
\uses{lem:affine-blowup-chart-evaluation,lem:affine-blowup-ideal-power-fraction}
\lean{Alterations.reesChartFractionEquiv}
\leanok
Evaluation gives an isomorphism of $R$-algebras
$\mathcal R(I)_{(gT)}\cong R[I/g]$.
\end{lemma}
\begin{proof}
Every element of the homogeneous localization has the form
$aT^n/(gT)^n$ with $a\in I^n$. The preceding power-fraction lemma shows that
its image is in $R[I/g]$. Conversely the image contains $R$, and contains
each generator $a/g$ as the image of $aT/(gT)$. Since it is a subalgebra, it
contains the algebra generated by these elements. Its image is therefore
exactly $R[I/g]$. The injectivity just proved identifies the rings.
\end{proof}

\endgroup

\begingroup

\begin{corollary}[The actual Proj chart]
\label{cor:affine-blowup-chart-scheme}
\uses{def:affine-blowup,lem:affine-blowup-chart-ring}
\lean{Alterations.affineBlowupChartIso}
\leanok
The basic open $D_+(gT)$ of $\operatorname{Bl}_I(\operatorname{Spec}R)$ is
isomorphic to $\operatorname{Spec}R[I/g]$.
\end{corollary}
\begin{proof}
The affine-open description of the projective spectrum identifies $D_+(gT)$
with $\operatorname{Spec}\mathcal R(I)_{(gT)}$. Compose this scheme isomorphism
with the spectrum of the inverse of the preceding ring isomorphism.
\end{proof}

\endgroup

\begingroup

\begin{lemma}[The chart lies over the original affine scheme]
\label{lem:affine-blowup-chart-over-base}
\uses{cor:affine-blowup-chart-scheme,lem:affine-blowup-chart-ring}
\lean{Alterations.affineBlowupChartInclusion_isOpenImmersion,Alterations.affineBlowupChartInclusion_structure}
\leanok
The map $\operatorname{Spec}R[I/g]\to\operatorname{Bl}_I(\operatorname{Spec}R)$ obtained from
the chart identification is an open immersion. Its composition with the blowup's
structure map is the morphism induced by the original inclusion $R\to R[I/g]$.
\end{lemma}
\begin{proof}
The chart map is the composition of the spectrum of a ring isomorphism with
the inclusion of the basic open $D_+(gT)$, so it is an open immersion. On
degree-zero elements the ring isomorphism is evaluation of constants, which
sends $a$ to the original coefficient $a\in R[I/g]$. The structure map of
the projective spectrum is induced on this chart by these degree-zero elements.
Thus the composite ring map from $R$ is exactly the coefficient inclusion,
and contravariance of the spectrum gives the stated equality of morphisms.
\end{proof}
\endgroup

\begingroup

The blowup used in Section~4.8 of \cite{deJong1996} is a modification.
The argument below proves dominance directly for the actual affine blowup
of a nonzero ideal in a domain, and proves surjectivity when the ideal is
finitely generated.

\begin{lemma}[The chart is dominant over the base]
\label{lem:affine-blowup-chart-dominant}
\uses{lem:affine-blowup-chart-coefficients-injective,lem:affine-blowup-chart-over-base}
\lean{Alterations.affineBlowupChartInclusion_comp_isDominant}
\leanok
If $0\ne g\in I$, then the composition
$\operatorname{Spec}R[I/g]\to\operatorname{Bl}_I(\operatorname{Spec}R)
\to\operatorname{Spec}R$ is dominant.
\end{lemma}
\begin{proof}
The chart's compatibility with the structure map identifies the composition
with the spectrum map of the coefficient inclusion $R\to R[I/g]$.
For any ring map, the closure of its spectrum image is the zero locus of
its kernel. The preceding injectivity lemma makes this kernel zero, whose
zero locus is the entire spectrum. The composition thus has dense image.
\end{proof}

\endgroup

\begingroup

The blowup used in Section~4.8 of \cite{deJong1996} is a modification.
The argument below proves dominance directly for the actual affine blowup
of a nonzero ideal in a domain, and proves surjectivity when the ideal is
finitely generated.

\begin{theorem}[A nonzero ideal has dominant blowup]
\label{thm:affine-blowup-dominant}
\uses{lem:affine-blowup-chart-dominant}
\lean{Alterations.affineBlowupStructure_isDominant}
\leanok
Let $R$ be a domain and $I\ne0$ an ideal. The actual blowup structure map
$\operatorname{Bl}_I(\operatorname{Spec}R)\to\operatorname{Spec}R$ is dominant.
No finite generation assumption is required.
\end{theorem}
\begin{proof}
Since $I$ is nonzero, choose an element $g\in I$ with $g\ne0$.
The preceding lemma makes the composite through its actual chart dominant.
The image of this composite is contained in the image of the blowup
structure map. Taking closures shows that the latter image is dense as well.
\end{proof}

\endgroup

\begingroup

The blowing up of the ideal sheaf in \cite{deJong1996}, 4.8, is a
scheme over the original scheme, with its affine Rees constructions
identified on every affine restriction. We construct this scheme by
gluing the actual affine Rees projective spectra. Throughout, $X$ is
an arbitrary scheme and $\mathcal I$ is a native ideal sheaf on $X$.
No gluing datum or local comparison isomorphism is supplied as an input.

\begin{lemma}[Actual ideal transitions on affine opens {\cite[p.~67]{deJong}}]
\label{lem:blowup-gluing-ideal-transitions}
\lean{Alterations.Geometry.BlowupGluing.coefficientMap,Alterations.Geometry.BlowupGluing.center_map,Alterations.Geometry.BlowupGluing.coefficientMap_id,Alterations.Geometry.BlowupGluing.coefficientMap_comp}
\leanok
For an affine principal-open inclusion $U\subset V$, let
$r_{VU}:\Gamma(V,\mathcal O_X)\to\Gamma(U,\mathcal O_X)$ be the
actual restriction map. Then
$\mathcal I(V)\Gamma(U,\mathcal O_X)=\mathcal I(U)$.
The restriction maps preserve identities and compose in the prescribed order.
\end{lemma}
\begin{proof}
The compatibility of the native ideal sheaf with affine restriction
gives the ideal equality. The identity and composition equalities are
the corresponding equalities for the structure sheaf's restriction maps.
\end{proof}

\endgroup

\begingroup

Let $R$ be a commutative ring, $I\subset R$ an ideal, and $s\in R$.
Put $R_s=R[s^{-1}]$ and $I_s=IR_s$. No domain or finite-generation
hypothesis is imposed, and $s=0$ is allowed.

\begin{definition}[Coefficient map on Rees algebras]
\label{def:rees-localization-map}
\uses{lem:affine-rees-grading}
\lean{Alterations.Algebra.ReesLocalization.reesMap,Alterations.Algebra.ReesLocalization.reesGradedMap,Alterations.Algebra.ReesLocalization.reesLocalizationAlgebra}
\leanok
A ring homomorphism $f:R\to S$ induces
$\mathcal R_R(I)\to\mathcal R_S(IS)$ by coefficientwise polynomial
mapping. It preserves every homogeneous degree. For the localization
homomorphism $R\to R_s$, use this actual map to equip
$\mathcal R_{R_s}(I_s)$ with its $\mathcal R_R(I)$-algebra structure.
\end{definition}
\begin{proof}
The coefficient of degree $n$ of the image is $f(a_n)$, where
$a_n\in I^n$. Since extension of ideals commutes with powers,
$f(a_n)\in(IS)^n$. Hence the polynomial image belongs to the actual
target Rees subalgebra. Polynomial mapping is a ring homomorphism.
The image of a monomial $aT^n$ is $f(a)T^n$, which establishes
compatibility with the native homogeneous pieces.
\end{proof}

\endgroup

\begingroup

\begin{lemma}[The actual Rees map defines a native Proj map]
\label{lem:rees-coefficient-irrelevant}
\uses{def:rees-localization-map}
\lean{Alterations.reesGradedMap_irrelevant_le_map}
\leanok
For a ring homomorphism $f:R\to S$ and an ideal $I\subset R$, the
irrelevant ideal of $\mathcal R_S(IS)$ lies in the extension of the
irrelevant ideal of $\mathcal R_R(I)$ under the actual coefficient map.
\end{lemma}
\begin{proof}
It suffices to treat a homogeneous element $bT^n$ with $n>0$ and
$b\in(IS)^n=(I^n)S$. Membership in this extended ideal expresses $b$
as a finite sum of $S$-multiples of images of elements $a\in I^n$.
The corresponding $bT^n$ is the same sum of coefficient multiples of
the images of $aT^n$. Each $aT^n$ is a genuine positive homogeneous
Rees element, hence belongs to the original irrelevant ideal. Every
summand and their sum therefore belong to its extended ideal. Taking
all positive homogeneous components proves the containment.
\end{proof}

\endgroup

\begingroup

\begin{definition}[The actual induced blowup morphism]
\label{def:affine-blowup-coefficient-map}
\uses{lem:rees-coefficient-irrelevant}
\lean{Alterations.affineBlowupMap}
\leanok
Coefficient extension defines the native morphism
\[\operatorname{Bl}_{IS}(\operatorname{Spec}S)
  \longrightarrow\operatorname{Bl}_I(\operatorname{Spec}R)\]
induced by the actual graded coefficient map on Rees algebras.
\end{definition}

\endgroup

\begingroup

Let $A=\bigoplus_{n\geq0}A_n$ and $B=\bigoplus_{n\geq0}B_n$ be
commutative graded rings. A graded homomorphism $f:A\to B$ with
$B_+\subseteq f(A_+)B$ induces a morphism
$\operatorname{Proj}(B)\to\operatorname{Proj}(A)$. The statements below
identify its structural morphism to the spectrum of the degree-zero ring.
They supply compatibility over the coefficient ring when a projective
presentation is constructed from a homogeneous algebra presentation.

\begin{definition}[Restriction to degree zero]
\label{def:proj-degree-zero-map}
\lean{Alterations.Geometry.ProjStructureMap.degreeZeroMap}
\leanok
For a graded ring homomorphism $f:A\to B$, its degree-zero restriction
$f_0:A_0\to B_0$ sends an element $a$ to $f(a)$.
\end{definition}

\endgroup

\begingroup

Let $A=\bigoplus_{n\geq0}A_n$ and $B=\bigoplus_{n\geq0}B_n$ be
commutative graded rings. A graded homomorphism $f:A\to B$ with
$B_+\subseteq f(A_+)B$ induces a morphism
$\operatorname{Proj}(B)\to\operatorname{Proj}(A)$. The statements below
identify its structural morphism to the spectrum of the degree-zero ring.
They supply compatibility over the coefficient ring when a projective
presentation is constructed from a homogeneous algebra presentation.

\begin{lemma}[Constants in homogeneous localizations]
\label{lem:proj-localization-degree-zero}
\uses{def:proj-degree-zero-map}
\lean{Alterations.Geometry.ProjStructureMap.awayMap_comp_fromZero}
\leanok
For $s\in A$, let $f_s:(A_s)_0\to(B_{f(s)})_0$ be the homogeneous
localization homomorphism. If $\iota_A:A_0\to(A_s)_0$ and
$\iota_B:B_0\to(B_{f(s)})_0$ send a constant to its fraction with
denominator one, then $f_s\iota_A=\iota_B f_0$.
\end{lemma}
\begin{proof}
Both homomorphisms send $a\in A_0$ to the class of $f(a)/1$. Indeed the
homogeneous localization homomorphism sends the representative $a/1$ to
$f(a)/f(1)$, and $f(1)=1$. Equality on every element proves equality of
the ring homomorphisms.
\end{proof}

\endgroup

\begingroup

Let $A=\bigoplus_{n\geq0}A_n$ and $B=\bigoplus_{n\geq0}B_n$ be
commutative graded rings. A graded homomorphism $f:A\to B$ with
$B_+\subseteq f(A_+)B$ induces a morphism
$\operatorname{Proj}(B)\to\operatorname{Proj}(A)$. The statements below
identify its structural morphism to the spectrum of the degree-zero ring.
They supply compatibility over the coefficient ring when a projective
presentation is constructed from a homogeneous algebra presentation.

\begin{theorem}[Naturality of the projective structure morphism]
\label{thm:proj-structure-map-naturality}
\uses{lem:proj-localization-degree-zero}
\lean{Alterations.Geometry.ProjStructureMap.map_toSpecZero}
\leanok
Suppose $B_+\subseteq f(A_+)B$. Let
$\pi_A:\operatorname{Proj}(A)\to\operatorname{Spec}(A_0)$ and
$\pi_B:\operatorname{Proj}(B)\to\operatorname{Spec}(B_0)$ be the native
projective structure morphisms. Then
\[
\pi_A\circ\operatorname{Proj}(f)
=\operatorname{Spec}(f_0)\circ\pi_B.
\]
\end{theorem}
\begin{proof}
The opens $D_+(f(s))$ for positive-degree homogeneous $s\in A$ cover
$\operatorname{Proj}(B)$ by the irrelevant-ideal hypothesis. On this
cover the morphism $\operatorname{Proj}(f)$ is the spectrum of
$f_s:(A_s)_0\to(B_{f(s)})_0$, followed by the inclusion of $D_+(s)$.
The restriction of $\pi_A$ to $D_+(s)$ is $\operatorname{Spec}(\iota_A)$,
and the restriction of $\pi_B$ to $D_+(f(s))$ is
$\operatorname{Spec}(\iota_B)$. The preceding lemma and contravariant
functoriality of the spectrum identify the two composites on each member
of the cover. Morphisms of schemes agreeing after precomposition by every
member of an open cover are equal, proving the formula.
\end{proof}

\endgroup

\begingroup

\begin{lemma}[The genuine coefficient structure square]
\label{lem:affine-blowup-coefficient-square}
\uses{def:affine-blowup-coefficient-map,thm:proj-structure-map-naturality}
\lean{Alterations.reesGradedMap_degreeZero_constants,Alterations.affineBlowupMap_structure}
\leanok
The induced blowup morphism commutes with the native structure maps
to $\operatorname{Spec}R$ and $\operatorname{Spec}S$ and the spectrum
map of $f$. This concerns the specified coefficient map, not merely
an abstract morphism between the schemes.
\end{lemma}
\begin{proof}
The actual graded coefficient map sends the constant $r$ to the
constant $f(r)$. Restricting to degree zero thus commutes with the
canonical identifications of the two degree-zero Rees rings with
$R$ and $S$. Naturality of the native projective structure morphism,
followed by these coefficient identifications, gives the square.
\end{proof}

\endgroup

\begingroup

\begin{lemma}[Functoriality with explicitly named ideals]
\label{lem:affine-blowup-coefficient-functorial}
\uses{def:affine-blowup-coefficient-map,lem:affine-blowup-coefficient-square}
\lean{Alterations.reesGradedMapTo,Alterations.affineBlowupMapTo,Alterations.affineBlowupMapTo_structure,Alterations.affineBlowupMapTo_id,Alterations.affineBlowupMapTo_comp}
\leanok
If $f:R\to S$ carries $I$ to the extended ideal $J=IS$, denote the
induced native blowup map with both ideals specified. These specified
maps preserve identity and composition, and their structure squares
are the actual coefficient squares.
\end{lemma}
\begin{proof}
The coefficientwise polynomial map restricts to the genuine Rees algebra
of the named ideal $J$. Identity coefficients give the identity polynomial
map. For $R\xrightarrow{f}S\xrightarrow{g}T$, polynomial coefficient
mapping satisfies $g(f(P))=(g\circ f)(P)$. Restriction to the Rees
subalgebras therefore has these same identities, with the named source
and target ideals fixed. Native Proj carries these exact graded-map
identities to the contravariant identity and composition identities
of the scheme morphisms. The degree-zero coefficient calculation is
unchanged, so the specified maps have the same structure squares.
\end{proof}

\endgroup

\begingroup

The blowing up of the ideal sheaf in \cite{deJong1996}, 4.8, is a
scheme over the original scheme, with its affine Rees constructions
identified on every affine restriction. We construct this scheme by
gluing the actual affine Rees projective spectra. Throughout, $X$ is
an arbitrary scheme and $\mathcal I$ is a native ideal sheaf on $X$.
No gluing datum or local comparison isomorphism is supplied as an input.

\begin{definition}[The affine blowup diagram and its structure maps]
\label{def:blowup-gluing-affine-diagram}
\uses{lem:blowup-gluing-ideal-transitions,lem:affine-blowup-coefficient-functorial,lem:affine-blowup-coefficient-square}
\lean{Alterations.Geometry.BlowupGluing.affineDiagram,Alterations.Geometry.BlowupGluing.structureNatTrans}
\leanok
On the small affine Zariski site of $X$, assign to $U$ the actual
scheme $B_U=\operatorname{Proj}\mathcal R(\mathcal I(U))$.
To a principal-open inclusion $U\subset V$, assign the native
coefficient-induced morphism $B_U\to B_V$. These assignments define a
functor. Its native affine structure maps form a natural transformation
to the diagram $U\mapsto\operatorname{Spec}\Gamma(U,\mathcal O_X)$.
\end{definition}
\begin{proof}
The ideal-transition lemma makes each coefficient-induced morphism
well defined with the specified source and target. Identity and
composition of coefficient maps induce identity and composition of
the actual graded Rees maps and their native projective morphisms.
Thus the assignments form a functor. The coefficient structure square
for the native blowup morphism says precisely that its affine
structure maps give the stated natural transformation.
\end{proof}

\endgroup

\begingroup

The blowing up of the ideal sheaf in \cite{deJong1996}, 4.8, is a
scheme over the original scheme, with its affine Rees constructions
identified on every affine restriction. We construct this scheme by
gluing the actual affine Rees projective spectra. Throughout, $X$ is
an arbitrary scheme and $\mathcal I$ is a native ideal sheaf on $X$.
No gluing datum or local comparison isomorphism is supplied as an input.

\begin{lemma}[Actual affine-site restrictions are localizations]
\label{lem:blowup-gluing-coefficient-localization}
\lean{Alterations.Geometry.BlowupGluing.coefficientMap_isLocalization}
\leanok
For every arrow $U\subset V$ of the small affine Zariski site, there
is a section $s\in\Gamma(V,\mathcal O_X)$ such that its actual
restriction map realizes $\Gamma(U,\mathcal O_X)$ as ordinary
localization of $\Gamma(V,\mathcal O_X)$ at $s$.
\end{lemma}
\begin{proof}
By the definition of arrows in this site, $U=D(s)$ for such a
section $s$. The structure sheaf on an affine scheme identifies its
sections on $D(s)$ with localization at $s$, with exactly the actual
restriction map as coefficient map.
\end{proof}

\endgroup

\begingroup

This is the coefficient-localization compatibility needed to construct the
actual blowup used in \cite[4.8]{deJong}.

Let $R$ be any commutative ring, let $I\subset R$ be any ideal, and let
$s\in R$. Put $R_s=R[1/s]$ and $I_s=I\cdot R_s$. All constructions below
include zero divisors and $s=0$, and require no finite generation.

\begin{lemma}[The actual Rees coefficient map preserves the base]\label{lem:rees-localization-structure}
\uses{def:affine-blowup}
\lean{Alterations.Geometry.AffineBlowupBaseChange.reesGradedMap_degreeZero_constants,Alterations.Geometry.AffineBlowupBaseChange.reesGradedMap_projStructure}
\leanok
For every ring map $f:R\to S$, coefficientwise polynomial mapping gives
a graded map $\mathcal R(I)\to\mathcal R(I\cdot S)$. On degree zero it
is exactly $f$ under the actual coefficient-ring identifications.
Whenever the native Proj map is defined, its composite with the original
blowup structure map is the target blowup structure map followed by
$\operatorname{Spec}f$.
\end{lemma}
\begin{proof}
A constant polynomial $r$ is sent coefficientwise to the constant $f(r)$.
This proves the identity of actual degree-zero ring maps. Naturality of
the native projective structure map to the spectrum of the degree-zero
ring, followed by these constant-ring identifications, gives the stated
identity of scheme morphisms.
\end{proof}

\endgroup

\begingroup

This is the coefficient-localization compatibility needed to construct the
actual blowup used in \cite[4.8]{deJong}.

Let $R$ be any commutative ring, let $I\subset R$ be any ideal, and let
$s\in R$. Put $R_s=R[1/s]$ and $I_s=I\cdot R_s$. All constructions below
include zero divisors and $s=0$, and require no finite generation.

\begin{definition}[The actual localization comparison]\label{def:affine-blowup-localization-map}
\uses{lem:rees-localization-structure}
\lean{Alterations.Geometry.AffineBlowupBaseChange.affineBlowupLocalizationMap,Alterations.Geometry.AffineBlowupBaseChange.affineBlowupLocalizationMap_structure}
\leanok
The actual graded coefficient-localization map defines a canonical morphism
\[j:\operatorname{Bl}_{I_s}(\operatorname{Spec}R_s)
 \longrightarrow\operatorname{Bl}_I(\operatorname{Spec}R).\]
Writing $\pi_s$ and $\pi$ for the respective native structure maps, it
satisfies $\pi\circ j=\operatorname{Spec}(R\to R_s)\circ\pi_s$
as actual morphisms of schemes.
\end{definition}
\begin{proof}
The target Rees ring is the ordinary localization of the original Rees
ring at the degree-zero constant $s$. Every target positive homogeneous
element has a source homogeneous numerator after multiplication by a power
of this constant. Since this constant becomes a unit, the target irrelevant
ideal is contained in the extension of the source irrelevant ideal.
Thus the actual native Proj map exists. Its structure identity is the
preceding lemma for the original localization homomorphism.
\end{proof}

\endgroup

\begingroup

This is the coefficient-localization compatibility needed to construct the
actual blowup used in \cite[4.8]{deJong}.

Let $R$ be any commutative ring, let $I\subset R$ be any ideal, and let
$s\in R$. Put $R_s=R[1/s]$ and $I_s=I\cdot R_s$. All constructions below
include zero divisors and $s=0$, and require no finite generation.

\begin{theorem}[The exact localization image]\label{thm:affine-blowup-localization-open}
\uses{def:affine-blowup-localization-map}
\lean{Alterations.Geometry.AffineBlowupBaseChange.affineBlowupLocalizationMap_isOpenImmersion,Alterations.Geometry.AffineBlowupBaseChange.affineBlowupLocalizationMap_opensRange}
\leanok
The actual morphism $j$ is an open immersion with image exactly
$\pi^{-1}(D(s))$.
\end{theorem}
\begin{proof}
Ordinary localization at a degree-zero element gives an open immersion
on native Proj. On each genuine homogeneous affine chart the induced
homomorphism is localization at the same coefficient, so its spectrum
map is an open immersion; these charts cover Proj. On points, contraction
of homogeneous primes under localization is injective and has image the
homogeneous primes not containing the constant $s$. Hence the image is
$D_+(s)$. The native structure morphism pulls $D(s)$ back to this exact
basic open, proving the asserted formula.
\end{proof}

\endgroup

\begingroup

This is the coefficient-localization compatibility needed to construct the
actual blowup used in \cite[4.8]{deJong}.

Let $R$ be any commutative ring, let $I\subset R$ be any ideal, and let
$s\in R$. Put $R_s=R[1/s]$ and $I_s=I\cdot R_s$. All constructions below
include zero divisors and $s=0$, and require no finite generation.

\begin{theorem}[The actual restriction isomorphism over the localized base]\label{thm:affine-blowup-localization-iso}
\uses{thm:affine-blowup-localization-open}
\lean{Alterations.Geometry.AffineBlowupBaseChange.affineBlowupLocalizationIso,Alterations.Geometry.AffineBlowupBaseChange.affineBlowupLocalizationIso_hom_ι,Alterations.Geometry.AffineBlowupBaseChange.affineBlowupLocalizationIso_hom_structure}
\leanok
There is a canonical isomorphism
\[\operatorname{Bl}_{I_s}(\operatorname{Spec}R_s)\simeq\pi^{-1}(D(s))\]
whose forward map, followed by the actual open inclusion, is $j$.
After identifying $D(s)\simeq\operatorname{Spec}R_s$, this isomorphism
preserves the actual structure morphisms to $\operatorname{Spec}R_s$.
\end{theorem}
\begin{proof}
An open immersion identifies its source with its actual open image. Apply
this to $j$ and use the exact image formula above to obtain the isomorphism.
Its inclusion identity is the defining identity of the open-image
isomorphism. To check the structure maps to the localized base, compose
both with $\operatorname{Spec}R_s\to\operatorname{Spec}R$. This is the
canonical principal-open inclusion under $D(s)\simeq\operatorname{Spec}R_s$
and hence is a monomorphism. The two composites agree by the inclusion
identity, the restriction pullback identity, and the structure identity
for $j$. Cancelling that monomorphism proves compatibility over
$\operatorname{Spec}R_s$.
\end{proof}

\endgroup

\begingroup

This is the coefficient-localization compatibility needed to construct the
actual blowup used in \cite[4.8]{deJong}.

Let $R$ be any commutative ring, let $I\subset R$ be any ideal, and let
$s\in R$. Put $R_s=R[1/s]$ and $I_s=I\cdot R_s$. All constructions below
include zero divisors and $s=0$, and require no finite generation.

\begin{corollary}[The genuine base-change square]\label{cor:affine-blowup-localization-pullback}
\uses{thm:affine-blowup-localization-iso}
\lean{Alterations.Geometry.AffineBlowupBaseChange.isPullback_affineBlowupLocalization}
\leanok
The actual square with upper map $j$, left vertical map $\pi_s$, right
vertical map $\pi$, and lower map $\operatorname{Spec}(R\to R_s)$ is Cartesian.
\end{corollary}
\begin{proof}
Restriction to $D(s)$ is the defining pullback of $\pi$ along its open
inclusion. Replace the upper-left corner by the preceding actual blowup
isomorphism and the lower-left corner by the canonical isomorphism
$D(s)\simeq\operatorname{Spec}R_s$. The two compatibility identities from
the theorem identify the transported arrows with exactly $j$ and $\pi_s$.
Pullback squares are invariant under these corner isomorphisms, proving
the assertion for the native morphisms.
\end{proof}

\endgroup

\begingroup

Let $R\to S$ be an algebra of commutative rings, let $s\in R$,
and suppose that the actual structure map exhibits $S$ as a
localization of $R$ at the powers of $s$. Let $I$ be an
arbitrary ideal of $R$. No domain, nonzero, or finite generation
hypotheses are imposed.

\begin{definition}[The actual Rees coefficient algebra]
\label{def:rees-coefficient-localization-algebra}
\uses{def:rees-localization-map}
\lean{Alterations.Algebra.ReesCoefficientLocalization.reesCoefficientAlgebra}
\leanok
Equip $\mathcal R_S(IS)$ with its $\mathcal R_R(I)$-algebra
structure induced by the actual coefficientwise polynomial map
$\mathcal R_R(I)\to\mathcal R_S(IS)$.
\end{definition}

\endgroup

\begingroup

Let $R\to S$ be an algebra of commutative rings, let $s\in R$,
and suppose that the actual structure map exhibits $S$ as a
localization of $R$ at the powers of $s$. Let $I$ be an
arbitrary ideal of $R$. No domain, nonzero, or finite generation
hypotheses are imposed.

\begin{lemma}[Actual Rees numerators over an arbitrary localization ring]
\label{lem:rees-coefficient-localization-denominators}
\uses{def:rees-coefficient-localization-algebra}
\lean{Alterations.Algebra.ReesCoefficientLocalization.reesCoefficient_surj}
\leanok
For every $p\in\mathcal R_S(IS)$ there are $N\ge0$ and
$q\in\mathcal R_R(I)$ with $p\,s^N=f(q)$.
\end{lemma}
\begin{proof}
The coefficient $b_n$ of $p$ lies in $(IS)^n=(I^n)S$.
The extended-ideal membership criterion for the actual localization
$R\to S$ gives $a_n\in I^n$ and $k_n\ge0$ with
$b_ns^{k_n}=f(a_n)$. Thus the degree-$n$ component multiplied
by $s^{k_n}$ is the image of the actual Rees monomial $a_nT^n$.
There are only finitely many nonzero polynomial components.
Denominators of a finite sum can be combined using
\[
 (x+y)s^{k+\ell}=f(as^\ell+bs^k)
 \quad\text{if }xs^k=f(a),\ ys^\ell=f(b).
\]
The zero sum uses numerator zero and exponent zero. Induction on
the actual finite polynomial support proves the assertion.
\end{proof}

\endgroup

\begingroup

Let $R\to S$ be an algebra of commutative rings, let $s\in R$,
and suppose that the actual structure map exhibits $S$ as a
localization of $R$ at the powers of $s$. Let $I$ be an
arbitrary ideal of $R$. No domain, nonzero, or finite generation
hypotheses are imposed.

\begin{lemma}[Actual Rees equality over an arbitrary localization ring]
\label{lem:rees-coefficient-localization-equality}
\uses{def:rees-coefficient-localization-algebra}
\lean{Alterations.Algebra.ReesCoefficientLocalization.reesCoefficient_exists_of_eq}
\leanok
If $q,q'\in\mathcal R_R(I)$ have equal images in
$\mathcal R_S(IS)$, then $s^Nq=s^Nq'$ for some $N\ge0$.
\end{lemma}
\begin{proof}
The polynomial algebra $S[T]$ is the actual coefficient
localization of $R[T]$ at the constant polynomial $s$.
Equality of the images of $q,q'$ in its Rees subalgebra
therefore gives $s^Nq=s^Nq'$ in $R[T]$ by the ordinary
localization equality criterion. The inclusion of the actual
source Rees subalgebra into $R[T]$ is injective, so this is
also equality of the original Rees elements.
\end{proof}

\endgroup

\begingroup

Let $R\to S$ be an algebra of commutative rings, let $s\in R$,
and suppose that the actual structure map exhibits $S$ as a
localization of $R$ at the powers of $s$. Let $I$ be an
arbitrary ideal of $R$. No domain, nonzero, or finite generation
hypotheses are imposed.

\begin{theorem}[Native Rees localization over any coefficient model]
\label{thm:rees-coefficient-localization}
\uses{lem:rees-coefficient-localization-denominators,lem:rees-coefficient-localization-equality}
\lean{Alterations.Algebra.ReesCoefficientLocalization.reesCoefficient_isLocalization}
\leanok
With the actual coefficient algebra structure,
$\mathcal R_S(IS)$ is the native localization of
$\mathcal R_R(I)$ at its degree-zero constant $s$.
\end{theorem}
\begin{proof}
The image of $s$ is a unit in $S$, so its further image in
the target Rees ring is a unit. The denominator lemma supplies
a source numerator for each target element after multiplication
by a power of $s$. The equality lemma supplies the power-of-$s$
equality criterion. These are the three defining conditions
for the native localization structure. The proof uses the
given coefficient algebra map throughout and requires no
chosen isomorphism between $S$ and a canonical model.
\end{proof}
\endgroup

\begingroup

This is the coefficient-localization compatibility needed to construct the
actual blowup used in \cite[4.8]{deJong}.

Let $R$ be any commutative ring, let $I\subset R$ be any ideal, and let
$s\in R$. Put $R_s=R[1/s]$ and $I_s=I\cdot R_s$. All constructions below
include zero divisors and $s=0$, and require no finite generation.

\begin{theorem}[Independence of the localization presentation]\label{thm:affine-blowup-general-localization}
\uses{thm:affine-blowup-localization-open,cor:affine-blowup-localization-pullback,thm:rees-coefficient-localization}
\lean{Alterations.Geometry.AffineBlowupBaseChange.affineBlowupMap_isOpenImmersion_of_isLocalization,Alterations.Geometry.AffineBlowupBaseChange.affineBlowupMap_opensRange_of_isLocalization,Alterations.Geometry.AffineBlowupBaseChange.isPullback_affineBlowupMap_of_isLocalization,Alterations.Geometry.AffineBlowupBaseChange.isPullback_affineBlowupMapTo_of_isLocalization}
\leanok
Let $S$ be any $R$-algebra which is a localization of $R$ at the powers of
$s$, and put $J=I\cdot S$. The actual coefficient-induced morphism
$\operatorname{Bl}_J(\operatorname{Spec}S)\to\operatorname{Bl}_I(\operatorname{Spec}R)$
is an open immersion with image $\pi^{-1}(D(s))$. Its actual coefficient
structure square over $\operatorname{Spec}S\to\operatorname{Spec}R$ is Cartesian.
In particular the statement applies directly to the section ring of a
principal affine open, with its original restriction homomorphism. The same
Cartesian statement holds when the extended ideal is named $J$ and identified
by the equality $I\cdot S=J$, with the actual induced map to that named ideal.
\end{theorem}
\begin{proof}
The coefficientwise map on Rees rings is ordinary localization at the
constant $s$: clear a power of $s$ in each homogeneous coefficient and
then use a common power for the finitely many coefficients. Equality is
likewise detected after multiplication by a power of $s$; this remains
valid with zero divisors. Thus the actual graded map satisfies the general
Proj localization theorem, giving the asserted open immersion and exact
image. The spectrum map $\operatorname{Spec}S\to\operatorname{Spec}R$ is
an open immersion with image $D(s)$. The coefficient structure identity
makes the square commute, and its upper open image is exactly the inverse
image of the lower open image. The universal property of open subschemes
therefore identifies this actual square with the restriction pullback
square, proving that it is Cartesian. Substituting the actual equality of
center ideals gives the version with a separately named ideal $J$.
\end{proof}
\endgroup

\begingroup

The blowing up of the ideal sheaf in \cite{deJong1996}, 4.8, is a
scheme over the original scheme, with its affine Rees constructions
identified on every affine restriction. We construct this scheme by
gluing the actual affine Rees projective spectra. Throughout, $X$ is
an arbitrary scheme and $\mathcal I$ is a native ideal sheaf on $X$.
No gluing datum or local comparison isomorphism is supplied as an input.

\begin{lemma}[The affine blowup diagram is equifibered]
\label{lem:blowup-gluing-equifibered}
\uses{def:blowup-gluing-affine-diagram,lem:blowup-gluing-coefficient-localization,thm:affine-blowup-general-localization}
\lean{Alterations.Geometry.BlowupGluing.structureNatTrans_equifibered}
\leanok
For every arrow $U\subset V$ of the small affine Zariski site, the
square of actual affine blowup maps and actual affine structure maps
is a pullback square.
\end{lemma}
\begin{proof}
Such an arrow is a principal-open inclusion: $U=D(s)\subset V$ for a
section $s\in\Gamma(V,\mathcal O_X)$. The actual restriction ring
$\Gamma(U,\mathcal O_X)$ is ordinary localization of
$\Gamma(V,\mathcal O_X)$ at $s$. The native ideal-transition equality
identifies its centre ideal with the extension of $\mathcal I(V)$.
The affine blowup coefficient-localization theorem therefore gives
the asserted pullback square with exactly these four native maps.
\end{proof}

\endgroup

\begingroup

The blowing up of the ideal sheaf in \cite{deJong1996}, 4.8, is a
scheme over the original scheme, with its affine Rees constructions
identified on every affine restriction. We construct this scheme by
gluing the actual affine Rees projective spectra. Throughout, $X$ is
an arbitrary scheme and $\mathcal I$ is a native ideal sheaf on $X$.
No gluing datum or local comparison isomorphism is supplied as an input.

\begin{definition}[The global native blowup {\cite[p.~67]{deJong}}]
\label{def:global-native-blowup}
\uses{def:blowup-gluing-affine-diagram,lem:blowup-gluing-equifibered}
\lean{Alterations.Geometry.BlowupGluing.relativeGluingData,Alterations.Geometry.BlowupGluing.blowup,Alterations.Geometry.BlowupGluing.blowupStructure}
\leanok
Use the canonical affine isomorphisms
$U\simeq\operatorname{Spec}\Gamma(U,\mathcal O_X)$ to view the affine
structure maps as maps $B_U\to U$. They define a native relative
gluing datum over the locally directed cover of $X$ by its small
affine Zariski site. Define $\operatorname{Bl}_{\mathcal I}X$ to be
its glued scheme, and let $\pi:\operatorname{Bl}_{\mathcal I}X\to X$
be its native structure morphism.
\end{definition}
\begin{proof}
The affine structure maps are natural by the affine-diagram lemma.
Their squares are pullbacks by the equifibered lemma. Composing with
the canonical affine isomorphisms preserves these pullback squares.
The cover of $X$ by these affine opens is locally directed: at a
point of the intersection of two affine opens, a common principal
affine neighbourhood is contained in both. These are exactly the
hypotheses of native relative gluing, which constructs the glued
scheme and its morphism to $X$.
\end{proof}

\endgroup

\begingroup

The blowing up of the ideal sheaf in \cite{deJong1996}, 4.8, is a
scheme over the original scheme, with its affine Rees constructions
identified on every affine restriction. We construct this scheme by
gluing the actual affine Rees projective spectra. Throughout, $X$ is
an arbitrary scheme and $\mathcal I$ is a native ideal sheaf on $X$.
No gluing datum or local comparison isomorphism is supplied as an input.

\begin{theorem}[The global blowup has the prescribed affine pullbacks]
\label{thm:global-native-blowup-affine-pullback}
\uses{def:global-native-blowup}
\lean{Alterations.Geometry.BlowupGluing.blowupCover,Alterations.Geometry.BlowupGluing.chartInclusion,Alterations.Geometry.BlowupGluing.chartInclusion_isOpenImmersion,Alterations.Geometry.BlowupGluing.chartInclusion_structure,Alterations.Geometry.BlowupGluing.isPullback_chartOpen,Alterations.Geometry.BlowupGluing.isPullback_chart,Alterations.Geometry.BlowupGluing.chartInclusion_opensRange}
\leanok
For every affine open $U\subset X$, the native chart inclusion
$j_U:B_U\to\operatorname{Bl}_{\mathcal I}X$ is an open immersion,
and the square
\[
\begin{array}{ccc}
B_U&\xrightarrow{j_U}&\operatorname{Bl}_{\mathcal I}X\\
\downarrow&&\downarrow\pi\\
\operatorname{Spec}\Gamma(U,\mathcal O_X)&\longrightarrow&X
\end{array}
\]
is a pullback, where the bottom map is the canonical affine chart
map of $X$. The $j_U$ form an open cover and their images are exactly
$\pi^{-1}(U)$.
\end{theorem}
\begin{proof}
The native relative-gluing construction identifies each diagram
object with the inverse image of its corresponding base open.
Its canonical chart maps are open immersions and form an open cover
of the glued scheme. The relative-gluing structure equation gives
$j_U\pi=(B_U\to U)(U\hookrightarrow X)$, and its native pullback
theorem identifies this square as a pullback over $U$. Transporting
the lower-left vertex along the canonical affine isomorphism
$U\simeq\operatorname{Spec}\Gamma(U,\mathcal O_X)$ gives the displayed
square with the actual affine Rees structure map. The inverse-image
identity supplied by relative gluing gives the asserted exact image.
\end{proof}
\endgroup

\begingroup

Let $X$ be an integral scheme, let $D$ be native ideal-sheaf
data on $X$, and suppose that its closed support is not all
of $X$. Write $B=\operatorname{Bl}_D(X)$ for the actual
relative gluing of the affine Rees projective spectra and
$b:B\to X$ for its structure morphism. Finite generation is
not assumed.

\begin{lemma}[Dominance of each nonempty chart over the base]
\label{lem:global-blowup-integral-chart-dominant}
\uses{lem:integral-ideal-affine-nonzero,thm:affine-blowup-dominant,thm:global-native-blowup-affine-pullback}
\lean{Alterations.Geometry.BlowupIntegral.chartInclusion_comp_isDominant}
\leanok
For a nonempty affine open $U$, the actual chart inclusion
$j_U:B_U\to B$ has dominant composite $b j_U:B_U\to X$.
\end{lemma}
\begin{proof}
The native affine blowup structure map $B_U\to
\operatorname{Spec}\mathcal O_X(U)$ is dominant because
$D(U)$ is a nonzero ideal in a domain. The actual map from
that spectrum into $X$ has image $U$. A nonempty open in
an irreducible space is dense, so this second map is
dominant. Their composite is dominant. The actual global
structure-map triangle identifies it with $b j_U$.
\end{proof}

\endgroup

\begingroup

Let $X$ be an integral scheme, let $D$ be native ideal-sheaf
data on $X$, and suppose that its closed support is not all
of $X$. Write $B=\operatorname{Bl}_D(X)$ for the actual
relative gluing of the affine Rees projective spectra and
$b:B\to X$ for its structure morphism. Finite generation is
not assumed.

\begin{lemma}[Intersecting nonempty affine blowup charts]
\label{lem:global-blowup-integral-chart-intersections}
\uses{lem:global-blowup-integral-chart-dominant,thm:global-native-blowup-affine-pullback}
\lean{Alterations.Geometry.BlowupIntegral.chartInclusion_not_disjoint}
\leanok
For nonempty affine opens $U,V\subset X$, the open images
of $j_U$ and $j_V$ in $B$ intersect.
\end{lemma}
\begin{proof}
The dominant map $b j_U$ has image meeting the nonempty
open $V$. Choose $p\in B_U$ whose image in $X$ lies in
$V$. Then $j_U(p)$ belongs to the image of $j_U$ and
to $b^{-1}(V)$. The exact affine-chart pullback and image
formula identify $b^{-1}(V)$ with the open image of
$j_V$. This gives an actual point of the intersection.
\end{proof}

\endgroup

\begingroup

Let $X$ be an integral scheme, let $D$ be native ideal-sheaf
data on $X$, and suppose that its closed support is not all
of $X$. Write $B=\operatorname{Bl}_D(X)$ for the actual
relative gluing of the affine Rees projective spectra and
$b:B\to X$ for its structure morphism. Finite generation is
not assumed.

\begin{definition}[The cover by nonempty base affine pieces]
\label{def:global-blowup-integral-nonempty-cover}
\uses{def:global-native-blowup,thm:global-native-blowup-affine-pullback}
\lean{Alterations.Geometry.BlowupIntegral.nonemptyAffineBlowupCover}
\leanok
Restrict the actual affine blowup cover to affine opens
$U\subset X$ that are nonempty.
\end{definition}
\begin{proof}
An arbitrary point of $B$ lies in an original affine blowup
chart $B_U$. Its image under $b$ lies in $U$ by the actual
chart image formula. Thus that particular $U$ is nonempty.
The restricted family therefore still covers every point
of $B$. All maps and objects are the same native affine
blowups and open immersions as in the original cover.
\end{proof}

\endgroup

\begingroup

Let $X$ be an integral scheme, let $D$ be native ideal-sheaf
data on $X$, and suppose that its closed support is not all
of $X$. Write $B=\operatorname{Bl}_D(X)$ for the actual
relative gluing of the affine Rees projective spectra and
$b:B\to X$ for its structure morphism. Finite generation is
not assumed.

\begin{theorem}[Integrality of the global native blowup]
\label{thm:global-native-blowup-integral}
\uses{lem:global-blowup-integral-affine-pieces,lem:global-blowup-integral-chart-intersections,def:global-blowup-integral-nonempty-cover}
\lean{Alterations.Geometry.BlowupIntegral.isIntegral_blowup}
\leanok
The actual global blowup $B=\operatorname{Bl}_D(X)$ is an
integral scheme.
\end{theorem}
\begin{proof}
Every member of the restricted cover is integral, hence
reduced and irreducible. Reducedness descends through an
open cover. Continuous images of irreducible spaces are
irreducible, so each chart image is preirreducible.
The chart images form an open cover and have pairwise
nonempty intersections; the open-cover criterion for
preirreducibility therefore makes $B$ preirreducible.

The integral scheme $X$ has a point and an affine
neighborhood $U$ of that point. This $U$ is nonempty,
so $B_U$ is integral and therefore nonempty. Its chart
inclusion supplies a point of $B$. Consequently $B$ is
irreducible and reduced, which is precisely integrality.
No geometry of $B$ was supplied as a hypothesis.
\end{proof}
\endgroup
